\begin{table}[htbp]
\centering
\caption{Event-study estimates of the BAME--White majority gap in mental well-being}
\label{tab:event_study_bame}
\small
\begin{tabular}{lcc}
\toprule
Period & Event-study coefficient & p-value \\
\midrule
2019 (base) & 0.00 &  \\
2016 & 0.33** & 0.01 \\
2017 & 0.25* & 0.05 \\
2018 & 0.25** & 0.03 \\
2020/04 & -0.19 & 0.18 \\
2020/05 & -0.01 & 0.94 \\
2020/06 & -0.25* & 0.08 \\
2020/07 & -0.09 & 0.52 \\
2020/09 & 0.18 & 0.21 \\
2020/11 & 0.04 & 0.77 \\
2021/01 & 0.21 & 0.17 \\
2021/03 & 0.17 & 0.26 \\
2021/09 & 0.09 & 0.54 \\
Joint pre-trend &  & 0.05 \\
\bottomrule
\end{tabular}

\begin{flushleft}
\footnotesize
Notes: The table reports event-study interaction coefficients for the BAME--White majority gap in mental well-being, with 2019 omitted as the baseline period. The outcome is the April-adjusted and inverted GHQ-12 Likert score, so higher values indicate better mental well-being. A negative coefficient indicates that BAME respondents experienced a relative deterioration in mental well-being compared with the White majority. Standard errors are clustered at the individual level. *, **, and *** denote significance at the 10\%, 5\%, and 1\% levels, respectively. The row ``Joint pre-trend'' reports the p-value from the joint test of whether all pre-pandemic interaction coefficients are equal to zero.
\end{flushleft}
\end{table}